\centerline{\bf \Huge Алгебраические преобразования}

\begin{flushright}
        \scriptsize{}
\end{flushright}

\problem{1}
Про натуральные числа $a, b, c$ и $d$ известно, что

$$
2 a-6 b-2 c+3 d=4,3 a+8 b-3 c-4 d=6
$$


Докажите, что $d$ делится на $b$.

\problem{2}
Для некоторых натуральных чисел $a, b, c, d$ выполняются равенства

$$
\frac{a}{c}=\frac{b}{d}=\frac{a b+1}{c d+1}
$$


Докажите, что $a=c$ и $b=d$.

\problem{3}
Докажите, что произведение любых четырёх последовательных целых чисел при увеличении на 1 даёт точный квадрат.

\problem{4}
На каждой из ста карточек записано по одному числу, отличному от нуля, так, что каждое число равно квадрату суммы всех остальных. Какие это числа?

\problem{5}
Для произвольных натуральных чисел $x, y, z$ положим

$$
f(x, y, z)=x^{y^z}-x^{z^y}+y^{z^x}-y^{x^z}+z^{x^y}-z^{y^x}
$$


Вычислите сумму чисел $f(a, b, c)$ по всем возможным различным тройкам натуральных чисел ( $a, b, c$ ), не превосходящих 2024.

Числа в тройке могут совпадать, а тройки, отличающиеся друг от друга порядком чисел, (например, $(1,2,3)$ и $(2,1,3))$ считаются различными.

\problem{6}
Перемножаются все выражения вида $\pm \sqrt{1} \pm \sqrt{2} \pm \ldots \pm \sqrt{99} \pm \sqrt{100}$ (при всевозможных комбинациях знаков). Докажите, что результат является целым числом.

\problem{7}
По кругу стоят $2^n$ хамелеонов красного и синего цвета. Каждую минуту все хамелеоны, у которых соседи разного цвета, одновременно перекрашиваются в другой цвет: синие - в красный, красные - в синий. Остальные хамелеоны цвета не меняют. Докажите, что когда-нибудь все хамелеоны одновременно вернут себе первоначальный цвет.